\documentclass[11pt,a4paper]{article}
\usepackage[margin=2.2cm]{geometry}
\usepackage{amsmath,amssymb}
\usepackage{graphicx}
\usepackage{booktabs}
\usepackage{xcolor}
\usepackage{hyperref}
\usepackage{caption}
\usepackage[section]{placeins}
\graphicspath{{figs/}}
\newcommand{\HH}{H/H_0}

\title{DeepONet and Taylor surrogates for the robustness of the SHiP muon shield:\\[0.3ex]
\large a comparison on problems with a known answer and on the shield}
\author{}
\date{\today}

\begin{document}
\maketitle

\begin{abstract}
Two surrogates of the per-muon hit probability were used for the robustness analysis of the SHiP
muon shield: a DeepONet ensemble and a Taylor surrogate, a second-order expansion in the design
parameters with explicit gradient and Hessian coefficients plus a residual switched off at the
nominal design. On four toy problems with exact derivatives, both converge to the true gradient
and Hessian as the number of simulated samples grows. The DeepONet reaches a moderate accuracy
with the fewest samples and then levels off; the unweighted Taylor surrogates keep improving and
end with better Hessians on three of the four problems, and the reduced (low-rank) one reaches a
given Hessian accuracy with the fewest samples. The weighted Taylor surrogate used on the shield
needs 2--12 times more samples than the unweighted one. On the shield both surrogates reach the
same robustness conclusions, checked with the simulator: the same worst directions ($\cos=0.99$),
the same worst cases per tolerance within the errors up to $\pm3\%$, the same tolerance priorities
and the same momentum picture. The DeepONet predicts magnitudes more accurately (one-at-a-time
effects to $4\%$ against $6\%$); each surrogate predicts the other's worst designs better than its
own.
\end{abstract}

\tableofcontents

%==============================================================================
\section{The two surrogates}
%==============================================================================

Both model the probability that a muon $x$ reaches the sensitive plane for a design $u$ (the 43
free parameters, $u_i=\pm1$ at $\pm10\%$), and both are separable, $\text{logit}=\sum_k b_k(u)\,t_k(x)$,
so that the expected hits $H(u)$ and their exact derivatives come from one pass over the muons.
\begin{itemize}
  \item \textbf{DeepONet}: $b(u)$ and $t(x)$ are two networks (two hidden layers of 128 units, 32
        features). It makes no assumption about how $u$ enters.
  \item \textbf{Taylor}: $b(u)=[1,\,u,\,\tfrac12(Q^\top u)^2,\,\gamma(u)\,r(u)]$: a constant, the full
        gradient, a curvature of rank $r$ in a learned basis $Q$, and a residual network $r(u)$
        gated by $\gamma(u)=\langle u_i^4\rangle/\rho^4$, which vanishes with its gradient and Hessian at
        $u=0$. The gradient and Hessian at the nominal design are therefore explicit coefficients.
        The \emph{reduced} variant restricts the linear term to the same subspace ($Q^\top u$); the
        \emph{weighted} variant lets the Taylor coefficients learn only from designs near the nominal
        one (weight $\exp(-\langle u_i^4\rangle^{1/2}/0.3^2)$ per row) and the residual from all.
\end{itemize}
On the shield, each surrogate is an average of four seeds trained on the same simulated data
($15\,181$ designs, $6.2\times10^{10}$ muons); the Taylor surrogate there is the weighted full one
(rank 8).

%==============================================================================
\section{Toy problems: convergence to the exact derivatives}
%==============================================================================

Four problems with a binary outcome per sample and known probabilities (beam, short column, a
43-parameter synthetic problem, two moons; 10--43 parameters): each surrogate is trained on
$10^5$--$10^8$ simulated samples (1000 per design, designs in the region of interest pulled towards
$u=0$), and its gradient and Hessian at $u=0$ are compared with the exact ones (relative errors;
median of 2--3 seeds). The networks are smaller than on the shield (64 hidden units, 16 features),
the same for all budgets.

\begin{figure}[h]
  \centering
  \includegraphics[width=\linewidth]{compare_toys.png}
  \caption{Relative error of the gradient (top) and the Hessian (bottom) at $u=0$ against the
  number of simulated samples, on each toy problem (lines: median over the seeds; dots: seeds).}
  \label{fig:toys}
\end{figure}

\begin{table}[h]
\centering
\caption{Simulated samples needed to reach a Hessian error of $0.5$ and a gradient error of $0.1$
(log-log interpolation of the medians; ``$>10^8$'': not reached), and the errors at $10^8$ samples.}
\label{tab:toys}
\small
\resizebox{\linewidth}{!}{\begin{tabular}{llcccc}
\toprule
problem & surrogate & samples, Hessian $0.5$ & samples, gradient $0.1$ & Hessian at $10^8$ & gradient at $10^8$ \\
\midrule
beam         & DeepONet         & $2.6\times10^6$ & $6.6\times10^5$ & 0.33 & 0.029 \\
             & Taylor, weighted & $7.1\times10^6$ & $1.8\times10^6$ & 0.26 & 0.036 \\
             & Taylor, full     & $2.5\times10^6$ & $1.1\times10^6$ & 0.14 & 0.020 \\
             & Taylor, reduced  & $1.0\times10^6$ & $9.7\times10^5$ & 0.10 & 0.015 \\
\midrule
short column & DeepONet         & $3.9\times10^5$ & $1.5\times10^6$ & 0.28 & 0.064 \\
             & Taylor, weighted & $4.1\times10^7$ & $5.1\times10^6$ & 0.39 & 0.076 \\
             & Taylor, full     & $3.4\times10^6$ & $2.3\times10^6$ & 0.45 & 0.043 \\
             & Taylor, reduced  & $4.9\times10^5$ & $2.0\times10^6$ & 0.29 & 0.042 \\
\midrule
synthetic    & DeepONet         & $1.1\times10^7$ & $>10^8$ & 0.39 & 0.14 \\
             & Taylor, weighted & $6.3\times10^7$ & $>10^8$ & 0.44 & 0.25 \\
             & Taylor, full     & $2.0\times10^7$ & $>10^8$ & 0.31 & 0.16 \\
             & Taylor, reduced  & $8.4\times10^6$ & $>10^8$ & 0.34 & 0.20 \\
\midrule
two moons    & DeepONet         & $9.0\times10^5$ & $>10^8$ & 0.18 & 0.13 \\
             & Taylor, weighted & $1.8\times10^6$ & $>10^8$ & 0.09 & 0.18 \\
             & Taylor, full     & $9.8\times10^5$ & $>10^8$ & 0.05 & 0.12 \\
             & Taylor, reduced  & $2.9\times10^5$ & $>10^8$ & 0.06 & 0.49 \\
\bottomrule
\end{tabular}}
\end{table}

\paragraph{Every surrogate converges.} The errors of the gradient and the Hessian fall steadily
with the number of simulated samples for all four surrogates on all four problems
(Fig.~\ref{fig:toys}).

\paragraph{Few samples: the DeepONet.} At $10^5$ samples the DeepONet has the smallest errors
almost everywhere: the Taylor surrogates have more coefficients to determine (a full
gradient and a curvature per muon feature) before they are accurate. It reaches a gradient error of
$0.1$ first wherever that is reached.

\paragraph{Many samples: the Taylor surrogates.} The DeepONet levels off (beam Hessian $0.44$ at
$10^7$, $0.33$ at $10^8$), while the Taylor surrogates keep improving: at $10^8$ samples the full
Taylor has the better Hessian on three problems (beam $0.14$ vs.\ $0.33$, synthetic $0.31$ vs.\
$0.39$, two moons $0.05$ vs.\ $0.18$) and the worse one on the short column ($0.45$ vs.\ $0.28$).

\paragraph{Fewest samples for a given Hessian: the reduced Taylor.} Where the response has
low-rank structure, restricting the expansion to a subspace pays: the reduced Taylor reaches a
Hessian error of $0.5$ with the fewest samples on three problems (beam $1.0\times10^6$ vs.\
$2.6\times10^6$ for the DeepONet; synthetic $8.4\times10^6$ vs.\ $1.1\times10^7$; two moons
$2.9\times10^5$ vs.\ $9.0\times10^5$) and as fast as the DeepONet on the fourth. Its gradient fails
where it does not lie in the subspace (two moons: $0.49$).

\paragraph{The weighting costs samples.} Letting the Taylor coefficients learn from the designs
near $u=0$ only discards most of the data for them: the weighted surrogate needs 2--12 times more
samples than the unweighted one for the same Hessian accuracy, and is still somewhat less accurate
at $10^8$ samples. On the shield, with $6\times10^{10}$ simulated muons, it was the better choice
near the nominal design (Sec.~\ref{sec:shield}).

%==============================================================================
\section{The muon shield}
\label{sec:shield}
%==============================================================================

\subsection{Accuracy}

\begin{table}[h]
\centering
\caption{Surrogate vs.\ simulator on the reference designs (never trained on): median
$|\log(\text{pred}/\text{sim})|$ (signed for the nominal design), and the relative vector errors of the
43 one-at-a-time (OAT) secants and curvatures on the check sample (Poisson noise: $1.6\%$ and
$7.6\%$).}
\label{tab:val}
\begin{tabular}{lcc}
\toprule
 & DeepONet & Taylor \\
\midrule
nominal design, nominal muons      & $-2.0\%$ & $+1.3\%$ \\
designs in the half box            & $4.1\%$  & $8.5\%$ \\
designs in the box                 & $12.6\%$ & $12.0\%$ \\
check sample: nominal design       & $+0.6\%$ & $-1.5\%$ \\
check sample: designs in the box   & $3.9\%$  & $6.0\%$ \\
OAT secants                        & $3.9\%$  & $6.4\%$ \\
OAT curvatures                     & $12\%$   & $15\%$ \\
members' gradients agree ($\cos$)  & 0.97--0.99 & 0.90--0.95 \\
\bottomrule
\end{tabular}
\end{table}

\begin{figure}[h]
  \centering
  \includegraphics[width=\linewidth]{refs_lcso_tr4n30r_oat.png}\\[1ex]
  \includegraphics[width=0.75\linewidth]{refs_lcso_tr4n30r_calibration.png}
  \caption{Top: every parameter alone at $\pm10\%$ on the check sample: simulator and both
  surrogates. Bottom: predicted vs.\ simulated hits of the reference designs.}
  \label{fig:val}
\end{figure}

Each robustness analysis simulated the designs its own surrogate pointed to (walks along its
directions, concrete worst cases, worst designs within each tolerance), on the same $10^9$ muons.
Both surrogates can be scored on both sets:

\begin{table}[h]
\centering
\caption{Median $|\log(\text{pred}/\text{sim})|$ of each surrogate on the designs simulated in each
analysis. The scans are the same designs in both.}
\label{tab:cross}
\resizebox{\linewidth}{!}{\begin{tabular}{lcccc}
\toprule
 & \multicolumn{2}{c}{designs of the DeepONet analysis} & \multicolumn{2}{c}{designs of the Taylor analysis} \\
 & DeepONet & Taylor & DeepONet & Taylor \\
\midrule
walks within one tolerance unit (28)    & $5.0\%$ & $4.0\%$ & $3.4\%$ & $4.5\%$ \\
worst direction scaled to 1--10\% (7)   & $9.4\%$ & $4.7\%$ & $5.4\%$ & $12\%$ \\
worst direction, top changes only (4)   & $11\%$  & $30\%$  & $11\%$  & $20\%$ \\
single parameters at $\pm10\%$ (8)      & $5.5\%$ & $3.8\%$ & $3.5\%$ & $8.2\%$ \\
random designs (24)                     & $9.1\%$ & $6.6\%$ & $7.0\%$ & $12\%$ \\
worst designs within $\pm1$--$3\%$ (3)  & $4.0\%$ & $11\%$  & $5.8\%$ & $2.0\%$ \\
scans (32)                              & $3.7\%$ & $5.1\%$ & $3.7\%$ & $5.1\%$ \\
\bottomrule
\end{tabular}}
\end{table}

Each surrogate predicts the \emph{other's} designs better than its own (Table~\ref{tab:cross}): the
designs an analysis chooses are those its surrogate rates worst, which selects where it errs (the
optimiser's curse). Averaged over both sets the two are comparable near the nominal design; the
DeepONet is better on corner-like designs (the truncated directions) and on the one-at-a-time
scans. Along its own harmful directions the Taylor surrogate is $13$--$20\%$ low at one tolerance
unit, the DeepONet within $5\%$.

\subsection{Robustness conclusions}

\begin{table}[h]
\centering
\caption{The robustness conclusions of the two analyses. Simulated values on $10^9$ muons;
the cosines compare the two surrogates' quantities.}
\label{tab:conclusions}
\small
\begin{tabular}{p{0.38\linewidth}p{0.26\linewidth}p{0.26\linewidth}}
\toprule
 & DeepONet & Taylor \\
\midrule
gradient at $u=0$ & \multicolumn{2}{c}{$\cos=0.92$ between the two} \\
worst direction / top Hessian eigenvector / active direction & \multicolumn{2}{c}{$\cos=0.99$ / $0.98$ / $1.00$} \\
Hessian eigenvalues ($H_0$ per unit$^2$) & 3.41, 1.01, 0.86 & 2.39, 1.40, 0.58 \\
most sensitive parameters (top 8, 7 shared; rank corr.\ 0.93) & M7 field, M2 entrance, M7 length, M7 inner yoke, M3 inner yoke, M7 exit, M2 field, M6 length & M7 field, M7 exit, M7 length, M2 entrance, M7 inner yoke, M2 field, M4 length, M3 inner yoke \\
worst direction, one tolerance unit (sim.) & $\times4.95\pm0.19$ & $\times5.20\pm0.20$ \\
corner pattern, every parameter $\pm1.5\%$ (sim.) & $\times2.01$ & $\times2.24$ \\
worst single parameter (sim.) & M7 field, $\times2.11$ & M7 field, $\times2.29$ \\
worst design within $\pm1/2/3\%$ (sim.) & $\times1.72$ / $3.78$ / $9.0$ & $\times1.67$ / $3.62$ / $8.9$ \\
worst design within $\pm5/10\%$ (sim.) & $\times132$ / $3440$ & $\times103$ / $5090$ \\
tolerance allocation, worst $\le3/2/1.5\,H_0$ (sim.) & $2.44$ / $1.56$ / $1.40$ & $2.25$ / $1.64$ / $1.38$ \\
tightest tolerances & M7 field, M2 entrance, M7 length, M7 inner yoke, M7 exit & M2 entrance, M7 field, M7 length, M7 inner yoke, M3 inner yoke \\
random $\pm10\%$: median, 99th percentile (surr.) & $\times2.08$, $\times29$ & $\times2.03$, $\times34$ \\
momentum: hits at $30$--$75$ / $>150$\,GeV (sim.) & $66\%$ / $24\%$ & $65\%$ / $24\%$ \\
intermediate vs.\ high momenta, worst directions & $\cos=-0.24$ to $-0.47$ & $\cos=-0.39$ to $-0.52$ \\
\bottomrule
\end{tabular}
\end{table}

\begin{figure}[h]
  \centering
  \includegraphics[width=0.49\linewidth]{tolerance_boxes_deeponet.png}\hfill
  \includegraphics[width=0.49\linewidth]{tolerance_boxes_taylor.png}
  \caption{Worst case for a given tolerance, found with each surrogate and simulated: DeepONet
  (left), Taylor (right).}
  \label{fig:tolerance}
\end{figure}

The two analyses agree on every conclusion (Table~\ref{tab:conclusions}, Fig.~\ref{fig:tolerance}):
\begin{itemize}
  \item the same worst direction, as harmful in the simulator ($\times5.0$ and $\times5.2$), and the
        same active direction (the lengths and fields of M1--M4, harmful on both sides);
  \item the same worst cases per tolerance within the errors up to $\pm3\%$; at $\pm5\%$ and
        $\pm10\%$ both searches find the same mechanism (narrower M7 yokes, wider M2 entrance and
        M3 yoke), at different corners of similar severity;
  \item the same tolerance priorities (M7 and M2's entrance first, the gaps between magnets last),
        and allocations that hold their targets in the simulator;
  \item the same momentum picture: the harmful directions act through the $30$--$75$\,GeV muons,
        against the high-momentum ones.
\end{itemize}
They differ in the details the simulator does not decide alone: the Hessian of the Taylor
surrogate is less dominated by one direction (eigenvalues $2.4$ and $1.4$ against $3.4$ and $1.0$),
and the order of the second-tier parameters.

%==============================================================================
\section{Conclusions}
%==============================================================================

\begin{enumerate}
  \item \textbf{Both surrogates converge to the exact gradient and Hessian} on problems with a
        known answer. The DeepONet needs the fewest samples for a moderate accuracy; the Taylor
        surrogates keep improving where it levels off, and the reduced one needs the fewest samples
        for a given Hessian accuracy when the response is low-rank.
  \item \textbf{Weighting the Taylor part towards the nominal design costs samples} (2--12 times
        on the toys) but, with the shield's large training set, gave the better local accuracy.
  \item \textbf{On the shield the conclusions do not depend on the surrogate}: worst directions,
        worst cases per tolerance, tolerance priorities and the momentum picture agree, all checked
        with the simulator.
  \item \textbf{The DeepONet is the more accurate predictor of magnitudes} (one-at-a-time effects,
        walks along the harmful directions); the Taylor surrogate offers explicit, interpretable
        derivatives at the nominal design at a modest cost in accuracy.
  \item \textbf{A surrogate is least reliable on the designs it selects}: each predicts the other's
        worst designs better than its own, so worst-case magnitudes must always be simulated.
\end{enumerate}

\end{document}
